% ===========================================================================
%  Préambule : extensions, couleurs et styles du rapport
% ===========================================================================

% ------------------------------------------------------------------ langue
\usepackage[T1]{fontenc}
\usepackage[utf8]{inputenc}
\usepackage[french]{babel}
\usepackage[autostyle=true]{csquotes}

% ----------------------------------------------------------------- polices
\usepackage{libertinus}
\usepackage[scaled=0.84]{sourcecodepro}
\usepackage{microtype}

% ------------------------------------------------------------- mise en page
\usepackage[a4paper, top=2.6cm, bottom=2.8cm, left=2.7cm, right=2.7cm,
            headheight=15pt, footskip=1.2cm]{geometry}
\usepackage{scrlayer-scrpage}
\usepackage{enumitem}
\usepackage{graphicx}
\usepackage{array, booktabs, tabularx, multirow}
\usepackage{caption, subcaption}
\usepackage{float}
\usepackage[locale=FR, group-minimum-digits=4]{siunitx}

% ----------------------------------------------------------------- couleurs
\usepackage[table]{xcolor}
\definecolor{marine}{HTML}{0B3954}    % titres, traits forts
\definecolor{ocean}{HTML}{1098AD}     % accent
\definecolor{ecume}{HTML}{E3FAFC}     % fonds clairs
\definecolor{corail}{HTML}{E8590C}    % mises en garde, touches
\definecolor{sable}{HTML}{FFF4E6}     % fond des encadrés d'alerte
\definecolor{ardoise}{HTML}{495057}   % texte secondaire
\definecolor{brume}{HTML}{F1F3F5}     % fond du code
\definecolor{vert}{HTML}{2F9E44}

% ------------------------------------------------------------------ dessins
\usepackage{tikz}
\usetikzlibrary{arrows.meta, positioning, shapes.geometric, shapes.multipart,
                fit, calc, backgrounds, matrix, babel}
\usepackage{pgfplots}
\usepackage{pgfplotstable}
\usepackage[edges]{forest}
\pgfplotsset{compat=1.17}

% -------------------------------------------------------------------- code
\usepackage{listings}
\lstdefinestyle{bataille}{
  basicstyle=\ttfamily\small,
  keywordstyle=\color{marine}\bfseries,
  commentstyle=\color{ardoise}\itshape,
  stringstyle=\color{corail},
  backgroundcolor=\color{brume},
  frame=leftline, framerule=2.5pt, rulecolor=\color{ocean},
  framesep=8pt, xleftmargin=14pt, xrightmargin=4pt,
  aboveskip=10pt, belowskip=8pt,
  showstringspaces=false, columns=fullflexible, keepspaces=true,
  breaklines=true, breakatwhitespace=true, tabsize=4,
  captionpos=b, numbers=none,
  literate={é}{{\'e}}1 {è}{{\`e}}1 {ê}{{\^e}}1 {à}{{\`a}}1 {ù}{{\`u}}1 {ç}{{\c c}}1
           {ô}{{\^o}}1 {î}{{\^i}}1 {û}{{\^u}}1 {É}{{\'E}}1 {È}{{\`E}}1 {À}{{\`A}}1
           {â}{{\^a}}1 {ë}{{\"e}}1 {ï}{{\"i}}1 {œ}{{\oe}}1 {Ç}{{\c C}}1
           {→}{{$\rightarrow$}}1 {×}{{$\times$}}1 {…}{{\dots}}1
}
\lstdefinelanguage{css}{
  morekeywords={button, window, label},
  morecomment=[s]{/*}{*/}, morestring=[b]", sensitive=true
}
\lstset{style=bataille, language=C,
  morekeywords={bool, uint8_t, int8_t, uint32_t, uint64_t, bs_board, bs_ship, bs_coord,
                bs_game, bs_status, bs_knowledge, bs_rng, bs_view, bs_controller,
                bs_difficulty, bs_player, bs_cell_look, bs_turn_report, gpointer,
                bs_phase, bs_orientation, bs_view_choice, cell_context,
                GtkWidget, gboolean}}
\renewcommand{\lstlistingname}{Extrait}
\renewcommand{\lstlistlistingname}{Table des extraits de code}
\newcommand{\code}[1]{\texttt{\small #1}}

% --------------------------------------------------------------- algorithmes
\usepackage[french, ruled, vlined, linesnumbered]{algorithm2e}
\SetKwComment{Com}{$\triangleright$\ }{}
\SetCommentSty{textit}
\SetAlFnt{\small}
\DontPrintSemicolon

% ------------------------------------------------------------------ encadrés
\usepackage[most]{tcolorbox}
\newtcolorbox{retenir}[1][À retenir]{
  enhanced, colback=ecume, colframe=ocean, coltitle=white,
  fonttitle=\sffamily\bfseries\small, title={#1},
  boxrule=0pt, leftrule=3pt, arc=1.5pt, left=8pt, right=8pt, top=6pt, bottom=6pt,
  attach boxed title to top left={yshift=-2mm, xshift=3mm},
  boxed title style={colback=ocean, arc=1.5pt, boxrule=0pt, left=4pt, right=4pt}
}
\newtcolorbox{attention}[1][Défaut de la version 1]{
  enhanced, colback=sable, colframe=corail, coltitle=white,
  fonttitle=\sffamily\bfseries\small, title={#1},
  boxrule=0pt, leftrule=3pt, arc=1.5pt, left=8pt, right=8pt, top=6pt, bottom=6pt,
  attach boxed title to top left={yshift=-2mm, xshift=3mm},
  boxed title style={colback=corail, arc=1.5pt, boxrule=0pt, left=4pt, right=4pt}
}

% ---------------------------------------------------------------- titres
\addtokomafont{disposition}{\color{marine}}
\setkomafont{chapter}{\Huge\sffamily\bfseries}
\setkomafont{section}{\Large\sffamily\bfseries}
\setkomafont{subsection}{\large\sffamily\bfseries}
\setkomafont{caption}{\small}
\setkomafont{captionlabel}{\small\sffamily\bfseries\color{marine}}
\renewcommand*{\chapterformat}{%
  \textcolor{ocean}{\thechapter}\enskip\textcolor{ocean!40}{\rule[-0.15em]{1.5pt}{1em}}\enskip}
\RedeclareSectionCommand[beforeskip=0pt, afterskip=1.4\baselineskip]{chapter}

% ------------------------------------------------------- en-têtes et pieds
\clearpairofpagestyles
\ihead{\small\sffamily\color{ardoise}Bataille Navale --- rapport de projet}
\ohead{\small\sffamily\color{ardoise}\leftmark}
\cfoot{\small\sffamily\color{ardoise}\pagemark}
\KOMAoptions{headsepline=0.4pt}
\setkomafont{headsepline}{\color{ocean!50}}
\automark[chapter]{chapter}
\renewcommand*{\chaptermarkformat}{}

% ------------------------------------------------------------------- listes
\setlist{itemsep=2pt, topsep=4pt}
\setlist[itemize,1]{label={\tikz[baseline=-0.1ex]\fill[ocean] (0,0) -- (0.85ex,0.55ex) -- (0,1.1ex) -- cycle;}}
\setlist[itemize,2]{label=\textcolor{ocean}{--}}

% ------------------------------------------------------------ bibliographie
% Le moteur BibTeX suffit pour une dizaine de références et évite de dépendre de biber.
\usepackage[backend=bibtex, style=numeric, sorting=none, maxbibnames=6]{biblatex}
\addbibresource{references.bib}

% -------------------------------------------------------------------- liens
\usepackage[colorlinks=true, linkcolor=marine, urlcolor=ocean, citecolor=ocean,
            pdftitle={Bataille Navale --- rapport de projet},
            pdfauthor={Dmytro Palahin, Ivanna Makam},
            pdfsubject={Jeu de bataille navale en C11 et GTK 3},
            pdfkeywords={C11, GTK 3, MVC, bataille navale, intégration continue}]{hyperref}
\usepackage[french, nameinlink, noabbrev]{cleveref}
% Les noms propres au document sont posés après ceux que cleveref installe pour le français.
\newcommand{\nomsdesrenvois}{%
  \crefname{lstlisting}{extrait}{extraits}\Crefname{lstlisting}{Extrait}{Extraits}%
  \crefname{listing}{extrait}{extraits}\Crefname{listing}{Extrait}{Extraits}%
  \crefname{algocf}{algorithme}{algorithmes}\Crefname{algocf}{Algorithme}{Algorithmes}}
\addto\extrasfrench{\nomsdesrenvois}
\AtBeginDocument{\nomsdesrenvois}

% Les flottants restent près du texte qui les cite.
\renewcommand{\topfraction}{0.9}
\renewcommand{\bottomfraction}{0.8}
\renewcommand{\textfraction}{0.08}
\renewcommand{\floatpagefraction}{0.8}

% ------------------------------------------------------- raccourcis divers
\newcommand{\vun}{\textsf{v1}}
\newcommand{\vdeux}{\textsf{v2}}
\newcommand{\fichier}[1]{\texttt{\small #1}}
\newcommand{\depot}{\url{https://github.com/DmytroPalahin/Battleship}}
\newcommand{\touche}[1]{\tikz[baseline=(k.base)]\node[draw=ardoise, rounded corners=2pt,
  inner sep=2.5pt, font=\sffamily\scriptsize, fill=brume] (k) {#1};}

% Styles communs des schémas TikZ.
\tikzset{
  >={Stealth[length=2.2mm]},
  boite/.style={draw=marine, line width=0.7pt, rounded corners=3pt, fill=white,
                align=center, font=\sffamily\small, inner sep=6pt},
  couche/.style={boite, fill=ecume},
  coeur/.style={boite, fill=marine, text=white, draw=marine},
  externe/.style={boite, draw=ardoise, dashed, text=ardoise},
  fleche/.style={->, line width=0.7pt, draw=marine},
  note/.style={font=\sffamily\footnotesize\itshape, text=ardoise, align=center},
}

% Marque d'une case exclue par déduction, dans les figures de plateau.
\newcommand{\drapeau}[2]{%
  \node[font=\sffamily\bfseries\small, text=corail] at (#1+0.5,#2+0.5) {$\times$};}

% Quadrillage d'un plateau de côté #1 pour les figures.
\newcommand{\plateau}[1]{%
  \draw[step=1, line width=0.3pt, draw=ardoise!55] (0,0) grid (#1,#1);
  \draw[line width=0.8pt, draw=marine] (0,0) rectangle (#1,#1);}
